\chapter{Unicidad y teoremas avanzados de límites}

\section{Unicidad de límites}

Una de las propiedades fundamentales de las sucesiones convergentes es que
convergen a un límite único. Esto puede parecer obvio a primera vista: ¿cómo
podría una sucesión acercarse arbitrariamente a dos números distintos? Sin
embargo, como ocurre con muchos hechos intuitivos en análisis, la prueba
rigurosa requiere un razonamiento cuidadoso con nuestras definiciones
épsilon-$N$.

La clave para demostrar la unicidad es la \emph{prueba por contradicción}.
Asumiremos que una sucesión converge a dos límites distintos $L$ y $M$, y
mostraremos que esto lleva a una imposibilidad. La estrategia consiste en
elegir épsilon igual a la mitad de la distancia entre $L$ y $M$, y luego
demostrar que la sucesión no puede permanecer tan cerca de ambos límites
simultáneamente.

\subsection{Nuevas herramientas}

\subsubsection{Prueba por contradicción: \texttt{by\_contra}}

La táctica \texttt{by\_contra} nos permite probar un enunciado asumiendo su
negación y derivando una contradicción. La sintaxis es
\mintinline{lean}{by_contra h}, que añade una hipótesis \texttt{h} con la
negación del objetivo actual y cambia el objetivo a \texttt{False}.

\subsubsection{Positividad del valor absoluto: \texttt{abs\_pos\_of\_nonzero}}

El teorema \texttt{abs\_pos\_of\_nonzero} establece que si $x \neq 0$,
entonces $0 < |x|$. Esto es esencial para trabajar con distancias entre
puntos distintos.

\subsection{Enunciado matemático}

\begin{theorem}[Unicidad del límite]
Si una sucesión $a : \mathbb{N} \to \mathbb{R}$ converge tanto a $L$ como a
$M$, entonces $L = M$.
\end{theorem}

\subsection{Demostración en Lean}

\begin{minted}[frame=single, linenos, fontsize=\small]{lean}
theorem LimUnique (a : ℕ → ℝ) (L M : ℝ)
    (aToL : SeqLim a L) (aToM : SeqLim a M) :
    L = M := by
  by_contra h
  have f0 : L - M ≠ 0 := by bound
  have f1 : 0 < |L - M| := by apply abs_pos_of_nonzero f0
  have f2 : 0 < |L - M| / 2 := by bound
  specialize aToL (|L - M| / 2) f2
  specialize aToM (|L - M| / 2) f2
  choose N1 hN1 using aToL
  choose N2 hN2 using aToM
  have f3 : N1 ≤ N1 + N2 := by bound
  have f4 : N2 ≤ N1 + N2 := by bound
  specialize hN1 (N1 + N2) f3
  specialize hN2 (N1 + N2) f4
  have f5 : |L - M| =
    |(L - a (N1 + N2)) + (a (N1 + N2) - M)| := by ring_nf
  have f6 : |(L - a (N1 + N2)) + (a (N1 + N2) - M)| ≤
    |(L - a (N1 + N2))| + |(a (N1 + N2) - M)| := by apply abs_add
  have f7 : |(L - a (N1 + N2))| =
    |-(a (N1 + N2) - L)| := by ring_nf
  have f8 : |-(a (N1 + N2) - L)| =
    |(a (N1 + N2) - L)| := by apply abs_neg
  linarith [f5, f6, f7, f8, hN1, hN2]
\end{minted}

\subsection{Prueba en lenguaje natural}

\begin{proof}
Supongamos, por contradicción, que $L \neq M$. Entonces $L - M \neq 0$, por
lo que $|L - M| > 0$, y por tanto $\varepsilon := |L - M|/2 > 0$.

Como $a(n) \to L$, existe $N_1$ tal que $|a(n) - L| < \varepsilon$ para todo
$n \geq N_1$. Como $a(n) \to M$, existe $N_2$ tal que $|a(n) - M| <
\varepsilon$ para todo $n \geq N_2$.

Sea $N = N_1 + N_2$. Para $n = N$ tenemos simultáneamente $|a(N) - L| <
\varepsilon$ y $|a(N) - M| < \varepsilon$. Observemos que:
\[
  |L - M| = |(L - a(N)) + (a(N) - M)|.
\]
Por la desigualdad triangular:
\[
  |L - M| \leq |L - a(N)| + |a(N) - M|
           = |a(N) - L| + |a(N) - M|
           < \varepsilon + \varepsilon = |L - M|.
\]
Esto nos da $|L - M| < |L - M|$, lo cual es una contradicción. Por tanto,
$L = M$.
\end{proof}

\subsection{Por qué importa la unicidad}

La unicidad de los límites es fundamental para todo el edificio del análisis.
Nos permite hablar de \emph{el} límite de una sucesión en lugar de \emph{un}
límite, y garantiza que nuestra noción de convergencia está bien definida. Sin
la unicidad, muchos teoremas y técnicas estándar fallarían o requerirían
modificaciones sustanciales.

\section{Eventualmente: Las sucesiones convergentes permanecen cerca de su límite}

Cuando una sucesión converge a un límite distinto de cero, no solo se acerca
arbitrariamente a ese límite, sino que \emph{eventualmente} también se aleja
de cero. Esta propiedad de «estar eventualmente acotada lejos de cero» es
crucial para muchos teoremas que involucran cocientes y recíprocos.

La intuición es sencilla: si una sucesión converge a algún valor no nulo $L$,
entonces en algún momento los términos de la sucesión deben tener valor
absoluto al menos la mitad del de $L$. No pueden simultáneamente acercarse a
$L$ y reducirse hacia cero.

\subsection{Enunciado matemático}

\begin{theorem}
Si $a : \mathbb{N} \to \mathbb{R}$ converge a $L$ con $L \neq 0$, entonces
existe $N$ tal que para todo $n \geq N$ se tiene $|a(n)| \geq |L|/2$.
\end{theorem}

\subsection{Enfoque estratégico}

La clave está en usar la condición de convergencia con $\varepsilon = |L|/2$.
Como $L \neq 0$, tenemos $|L| > 0$, por lo que este épsilon es positivo. La
condición de convergencia nos dice entonces que eventualmente
$|a(n) - L| < |L|/2$, lo que por la desigualdad triangular inversa implica
que $|a(n)|$ es al menos $|L|/2$.

\subsection{Demostración en Lean}

\begin{minted}[frame=single, linenos, fontsize=\small]{lean}
theorem EventuallyGeHalfLim (a : ℕ → ℝ) (L : ℝ)
    (aToL : SeqLim a L) (LneZero : L ≠ 0) :
    ∃ N, ∀ n ≥ N, |L| / 2 ≤ |a n| := by
  specialize aToL (|L| / 2)
  have : 0 < |L| := by apply abs_pos_of_nonzero LneZero
  have : 0 < |L| / 2 := by bound
  specialize aToL this
  choose N hN using aToL
  use N
  intro n hn
  specialize hN n hn
  have l1 : |L| = |a n + (L - a n)| := by ring_nf
  have l2 : |a n + (L - a n)| ≤ |a n| + |L - a n| := by
    apply abs_add
  have l3 : |L - a n| = |-(a n - L)| := by ring_nf
  have l4 : |-(a n - L)| = |(a n - L)| := by apply abs_neg
  linarith [l1, l2, l3, l4, hN]
\end{minted}

\subsection{Prueba en lenguaje natural}

\begin{proof}
Como $L \neq 0$, tenemos $|L| > 0$, y por tanto $\varepsilon := |L|/2 > 0$.

Por la convergencia de $a$ hacia $L$, existe $N$ tal que
$|a(n) - L| < |L|/2$ para todo $n \geq N$.

Para cualquier $n \geq N$ tenemos:
\[
  |L| = |a(n) + (L - a(n))| \leq |a(n)| + |L - a(n)|
      = |a(n)| + |a(n) - L| < |a(n)| + \frac{|L|}{2}.
\]
Reordenando:
\[
  |L| - \frac{|L|}{2} < |a(n)|,
\]
por lo que $|a(n)| > |L|/2$, lo que implica $|a(n)| \geq |L|/2$.
\end{proof}

\subsection{Aplicaciones}

Este resultado es esencial para demostrar que el recíproco de una sucesión
convergente (con límite no nulo) es a su vez convergente. Garantiza que los
denominadores no se aproximan a cero, lo que haría que los recíprocos
divergieran.

\section{Continuidad del valor absoluto: Sucesiones de valores absolutos}

La función valor absoluto se comporta de manera muy natural respecto a los
límites: si una sucesión converge, la sucesión de los valores absolutos
converge al valor absoluto del límite. Esto es una manifestación de la
continuidad de la función valor absoluto.

La propiedad clave que usaremos es que el valor absoluto es Lipschitz
continuo con constante~1, es decir, $\big||x| - |y|\big| \leq |x - y|$ para
todo $x, y \in \mathbb{R}$. Esta desigualdad captura la idea de que el valor
absoluto no aumenta las distancias.

\subsection{Nuevas herramientas}

\subsubsection{Propiedad de Lipschitz: \texttt{abs\_Lipschitz}}

El teorema \texttt{abs\_Lipschitz} establece que para cualesquiera $x, y
\in \mathbb{R}$ se tiene $\big||x| - |y|\big| \leq |x - y|$. A veces se
denomina \emph{desigualdad triangular inversa} para valores absolutos.

\subsection{Enunciado matemático}

\begin{theorem}
Si $a : \mathbb{N} \to \mathbb{R}$ converge a $L$, y $b : \mathbb{N} \to
\mathbb{R}$ se define por $b(n) = |a(n)|$ para todo $n$, entonces $b$
converge a $|L|$.
\end{theorem}

\subsection{Demostración en Lean}

\begin{minted}[frame=single, linenos, fontsize=\small]{lean}
theorem AbsLim (a : ℕ → ℝ) (L : ℝ) (aToL : SeqLim a L)
    (b : ℕ → ℝ) (bEqAbsa : ∀ n, b n = |a n|) :
    SeqLim b |L| := by
  intro ε hε
  specialize aToL ε hε
  choose N hN using aToL
  use N
  intro n hn
  specialize hN n hn
  specialize bEqAbsa n
  rewrite [bEqAbsa]
  have : |(|a n|) - (|L|)| ≤ |a n - L| := by apply abs_Lipschitz
  bound
\end{minted}

\subsection{Prueba en lenguaje natural}

\begin{proof}
Sea $\varepsilon > 0$ dado. Como $a(n) \to L$, existe $N$ tal que
$|a(n) - L| < \varepsilon$ para todo $n \geq N$.

Para cualquier $n \geq N$:
\[
  |b(n) - |L|| = \big||a(n)| - |L|\big| \leq |a(n) - L| < \varepsilon,
\]
donde hemos usado la propiedad de Lipschitz del valor absoluto. Por tanto
$b(n) \to |L|$.
\end{proof}

\subsection{Por qué importa este resultado}

Este teorema es un caso especial de un principio mucho más general: las
\emph{funciones continuas preservan los límites}. La función valor absoluto
es continua en todas partes, por lo que transforma sucesiones convergentes en
sucesiones convergentes. Este principio se extiende a todas las funciones
continuas y es fundamental en el análisis matemático.

\section{Recíprocos de sucesiones convergentes}

Uno de los teoremas de límites más importantes concierne a los recíprocos: si
una sucesión converge a un límite no nulo, entonces la sucesión de los
recíprocos converge al recíproco del límite. Este resultado es crucial para
demostrar teoremas sobre cocientes y funciones racionales.

La prueba combina varias técnicas que hemos desarrollado: demostrar que la
sucesión permanece acotada lejos de cero (de modo que los recíprocos no
divergen), manipular cuidadosamente expresiones algebraicas con fracciones, y
usar la desigualdad triangular para controlar los términos de error.

\subsection{Nuevas herramientas}

\subsubsection{Valor absoluto de cocientes: \texttt{abs\_div}}

Para cualesquiera $x, y \in \mathbb{R}$ (con $y \neq 0$) se tiene
$|x/y| = |x|/|y|$. Esto permite separar los valores absolutos a través de
una división.

\subsubsection{No nulo a partir de valor absoluto positivo: \texttt{nonzero\_of\_abs\_pos}}

Si $0 < |x|$, entonces $x \neq 0$. Esto es útil para verificar que la
división es válida.

\subsection{Enunciado matemático}

\begin{theorem}
Si $a : \mathbb{N} \to \mathbb{R}$ converge a $L$ con $L \neq 0$, y $b :
\mathbb{N} \to \mathbb{R}$ se define por $b(n) = 1/a(n)$, entonces $b$
converge a $1/L$.
\end{theorem}

\subsection{Enfoque estratégico}

El reto es que la expresión $|1/a(n) - 1/L|$ involucra recíprocos, que
pueden ser difíciles de acotar directamente. Los pasos clave son:

\begin{enumerate}
  \item Usar el teorema \texttt{EventuallyGeHalfLim} para garantizar que
    $|a(n)| \geq |L|/2$ eventualmente.
  \item Elegir épsilon cuidadosamente: usar $\varepsilon \cdot |L|^2/2$ al
    aplicar la convergencia de $a$ hacia $L$.
  \item Simplificar algebraicamente:
    $|1/a(n) - 1/L| = |(L - a(n))/(a(n) \cdot L)|$.
  \item Separar el valor absoluto usando \texttt{abs\_div}.
  \item Acotar el denominador usando nuestra cota inferior de $|a(n)|$.
  \item Usar desigualdades encadenadas para mostrar que el resultado es
    menor que $\varepsilon$.
\end{enumerate}

\subsection{Demostración en Lean}

\begin{minted}[frame=single, linenos, fontsize=\small]{lean}
theorem InvLim (a : ℕ → ℝ) (L : ℝ) (aToL : SeqLim a L)
    (LneZero : L ≠ 0) (b : ℕ → ℝ)
    (bEqInva : ∀ n, b n = 1 / a n) :
    SeqLim b (1 / L) := by
  choose NhalfL hNhalfL using
    EventuallyGeHalfLim a L aToL LneZero
  intro ε hε
  have : 0 < |L| := by apply abs_pos_of_nonzero LneZero
  specialize aToL (ε * |L| * |L| / 2) (by bound)
  choose Na hNa using aToL
  use Na + NhalfL
  intro n hn
  specialize bEqInva n
  rewrite [bEqInva]
  have hnHalfL : NhalfL ≤ n := by bound
  have hna    : Na ≤ n      := by bound
  specialize hNhalfL n hnHalfL
  specialize hNa     n hna
  have : 0 < |a n|  := by bound
  have : a n ≠ 0    := by apply nonzero_of_abs_pos this
  have l1 : |1 / a n - 1 / L| =
      |(L - a n) / (a n * L)| := by field_simp
  have l2 : |(L - a n) / (a n * L)| =
      |L - a n| / |a n * L| := by apply abs_div
  have l3 : |L - a n| / |a n * L| =
      |L - a n| / (|a n| * |L|) := by bound
  have l4 : |L - a n| = |-(a n - L)| := by ring_nf
  have l5 : |-(a n - L)| = |a n - L| := by apply abs_neg
  have l6 : |a n - L| / (|a n| * |L|) <
      (ε * |L| * |L| / 2) / (|a n| * |L|) := by
        field_simp; nlinarith
  have l10 : |L - a n| / (|a n| * |L|) =
      |-(a n - L)| / (|a n| * |L|) := by rw [l4]
  have l11 : |-(a n - L)| / (|a n| * |L|) =
      |a n - L| / (|a n| * |L|) := by rw [l5]
  have l13 : ε * |L| * |L| / 2 / (|a n| * |L|) =
      ε * |L| / 2 / |a n| := by field_simp
  have l14 : ε * |L| / 2 / |a n| ≤ ε := by
        field_simp; bound
  linarith [l1, l2, l3, l10, l11, l6, l13, l14]
\end{minted}

\subsection{Prueba en lenguaje natural}

\begin{proof}
Sea $\varepsilon > 0$ dado. Como $L \neq 0$, por el teorema
\texttt{EventuallyGeHalfLim}, existe $N_1$ tal que $|a(n)| \geq |L|/2$ para
todo $n \geq N_1$.

Como $|L| > 0$, tenemos $\varepsilon \cdot |L|^2/2 > 0$. Por la convergencia
de $a$ hacia $L$, existe $N_2$ tal que $|a(n) - L| < \varepsilon |L|^2/2$
para todo $n \geq N_2$.

Sea $N = N_1 + N_2$. Para cualquier $n \geq N$ se tiene $|a(n)| \geq |L|/2
> 0$, así que $a(n) \neq 0$ y $b(n) = 1/a(n)$ está bien definido. Entonces:
\begin{align*}
  |b(n) - 1/L| &= |1/a(n) - 1/L|
    = \left|\frac{L - a(n)}{a(n) \cdot L}\right|
    = \frac{|L - a(n)|}{|a(n)| \cdot |L|}
    = \frac{|a(n) - L|}{|a(n)| \cdot |L|} \\
    &< \frac{\varepsilon |L|^2/2}{|a(n)| \cdot |L|}
    \leq \frac{\varepsilon |L|^2/2}{(|L|/2) \cdot |L|}
    = \frac{\varepsilon |L|^2/2}{|L|^2/2}
    = \varepsilon.
\end{align*}
Por tanto $b(n) \to 1/L$.
\end{proof}

\subsection{Prueba conseguida}

Acabas de completar una de las pruebas más desafiantes del análisis
elemental. El teorema del límite del recíproco es un hito importante: has
demostrado que los recíprocos preservan la convergencia (cuando el límite es
no nulo).

% ----
\subsection{Lo que lograste}

Esta prueba requirió orquestar múltiples técnicas sofisticadas:

\begin{itemize}
  \item Usar \texttt{EventuallyGeHalfLim} para garantizar que los
    denominadores permanezcan acotados lejos de cero.
  \item Elegir un épsilon cuidadosamente calibrado ($\varepsilon \cdot
    |L|^2/2$) para que el álgebra funcione.
  \item Manipular expresiones fraccionarias complejas con denominador común.
  \item Aplicar \texttt{abs\_div} para separar los valores absolutos a
    través de la división.
  \item Encadenar una secuencia de desigualdades para alcanzar la cota
    final.
\end{itemize}

Cada paso se apoyó en los anteriores, mostrando cómo las pruebas matemáticas
se construyen a partir de bloques cuidadosamente ensamblados.

\subsection{Aplicaciones y extensiones}

Con este teorema en mano, ya disponemos de un conjunto completo de
herramientas para los límites de funciones racionales. Combinado con los
resultados anteriores sobre sumas y productos, ahora podemos demostrar:

\begin{quote}
  Si $a(n) \to L$ y $c(n) \to M$ con $M \neq 0$, entonces
  $a(n)/c(n) \to L/M$.
\end{quote}

La prueba es directa: primero se muestra que $1/c(n) \to 1/M$ usando el
teorema del recíproco que acabamos de probar, y luego se usa el teorema del
producto para mostrar que $a(n) \cdot (1/c(n)) \to L \cdot (1/M) = L/M$.

Esto completa la aritmética fundamental de los límites: sumas, productos y
cocientes. Estos son los bloques de construcción para analizar límites de
polinomios, funciones racionales y expresiones mucho más complejas a lo largo
del cálculo y el análisis.

\subsection{Dominio de la técnica}

El teorema del recíproco ilustra una lección crucial en la demostración
matemática: a veces el épsilon «correcto» no es la elección obvia. La
expresión $\varepsilon \cdot |L|^2/2$ puede parecer misteriosa al principio,
pero está precisamente diseñada para que las desigualdades finales resulten
correctas. Este tipo de pensamiento estratégico —trabajar hacia atrás desde
lo que necesitas para saber qué debes asumir— está en el corazón de la
resolución de problemas matemáticos.

Ya dominas las técnicas esenciales para demostrar teoremas de límites.
¡Bien hecho!

%%% Local Variables:
%%% mode: LaTeX
%%% TeX-master: "Analisis_en_Lean"
%%% End:
